\subsection{Description et indexation automatique} 

Quoi décrire ?

à rajouter pour faire la transition : %L'apport de l'IA (Vers la classification sémantique) : La détection brute des coupures ne suffit plus ; l'enjeu actuel est la Shot semantics (décrire le contenu d'un plan et sa typologie) (Arnold  Tilton, 2019).

La description automatique d’images peut être centrée sur de nombreux éléments au sein d'une \gls{sequence} de film : dialogues, ambiance sonore, visages, objets, spatialité des éléments, actions, émotions... Toutefois, la définition d’un processus au sein d’un logiciel de gestion des collections impose un cadrage défini des fonctionnalités. Le corpus patrimonial retenu pour notre étude, constitué de \gls{rush} muets des premiers temps, exclue d’emblée l’analyse automatique de caractéristiques sonores. L’analyse automatique se concentre donc exclusivement sur les caractéristiques visuelles de l’image. Deux approches coexistent alors, celle de l’identification des éléments dans l’image de manière précise, ou celle d’éventuellement générer des résumés automatiques sur ces séquences, à une échelle plus globale. Or, la seconde approche, fondée sur des grands modèles de langage, implique des architectures génératives couteuses et lourdes. Nous choisissons ainsi d’adopter la première, afin de garantir une indexation fiable et stable, plus légère, à travers la détection et l’identification d’objets sur les \gls{imgcle}préalablement extraites. Cette tâche peut être réalisée par des architectures de réseaux de neurones relativement légers, pouvant extraire des concepts directement indexables au sein du logiciel.  



\subsubsection{Outils utilisés}

à utiliser ? %Le rôle des modèles (CNN et Détecteurs) : Pour analyser le contenu de ces plans, on utilise l'apprentissage supervisé (Supervised Learning), où des algorithmes apprennent à partir de données labellisées par l'humain (DeLua, 2024). C'est le cas des réseaux de neurones convolutifs (CNN) comme Faster R-CNN ou YOLO (un détecteur de type Single Shot qui prédit et classifie "en une seule fois") pour la détection d'objets ou de visages (Arnold  Tilton, 2019 ; Robert, 2026).

à intégrer ? : %au lieu d'utiliser des larges modèles qui tentent de tout classifier, plutot utiliser des modèles "low parameter", c'est-à-dire des modèles 3D-CNN légers, avec moins de paramètres dédiés à la reconnaissance d'une seule activité précise. prenennt moins de mémoire vidéo. \footcite{pattichis2023}

à intégrer ? : %amener l'algorithme aux données : C'est la discussion centrale de votre partie. Les outils d'AVP commerciaux (souvent basés sur le Cloud) sont inutilisables pour les archives en raison des droits de propriété intellectuelle et de la sécurité des données (van Noord et al., 2021). La solution réside dans des architectures locales comme DANE (projet CLARIAH), basées sur un système décentralisé de workers (van Noord et al., 2021). Ce type d'infrastructure produit des métadonnées exploitables, mais aussi des vecteurs de caractéristiques (représentations abstraites) ouvrant la voie à la recherche par similarité visuelle ou sonore (van Noord et al., 2021).

%Rajouter des exemples de projets qui utilisent YOLO dans le domaine des archives ? 

Si pour la tâche de segmentation automatique, nous avons écarté les architectures de modèles fondés sur des réseaux de neurones, la reconnaissance d'objets dans l'image suppose une prédiction fondée sur un apprentissage machine profond des réseaux de neurones. En effet, la plupart des outils de reconnaissance d'objets dans l'image sont fondés sur des modèles de réseaux de neurones convolutifs (CNN) dont nous avons exposé le fonctionnement dans la section précédente.  

Avantage de placer un réseau de neurones qu'à cette étape-là du processus (cf article 2sds) 

Par exemple, l’outil Snoop développé par l’INA est fondé sur le CNN ResNet. (c’est un moteur de recherche : pas pertinent ?) 

Nous détaillerons dans cette section l’utilisation de l’architecture \gls{yolo}, une série d’algorithmes dédiés à la détection d’objets visuels dans les images et les flux vidéo, et fondé sur les réseaux neuronaux convolutifs, développé par Ultralytics en 2015. 
%source = wikipédia mais trouver mieux  
%Réexpliquer ce que permettent les CNN ? 
C’est l’un des outils de détection d’objets les plus documentés, y compris dans le domaine des sciences humaines et du patrimoine. Fondé sur la \gls{Vision par ordinateur} et entraîné grâce à des méthodes d'\gls{Apprentissage profond}, \gls{yolo} est en accès libre, et permet notamment la détection d'objets et la classification d'images, permettant ainsi de poser les bases d’une indexation \footcite{jocher2022}. Le modèle extrait dans un premier temps les caractéristiques visuelles de l'image qui lui est donnée, les agrège, afin de finalement proposer une détection d'objets \footcite{norindr2023}. Les versions récentes du modèle telles que YoLo12 ne sont désormais plus seulement basées sur des \gls{cnn} mais sur des mécanismes d'auto-attention, similaires à ceux des \gls{Transformers} dont nous avions évoqué l'architecture \footcite{2025b}. Toutefois, comme nous l'avions évoqué, cette architecture demande une consommation de mémoire élevée et une puissance de calcul lourde \footcite{2025b}. Ultralytics, créateur du modèle, recommande des versions moins avancées pour une mise en application rapide, telles que la version Yolov8, uniquement fondée sur un réseau de neurones convolutif \footcite{2023c}. Yolov8 propose une version optimisée de cette architecture en trois temps : extraction, aggrégation, et détection \footcite{yaseen2024}. Le modèle se caractérise par sa capacité à analyser l'image en une seule fois, contrairement à d'autres modèles dont les résultats impliquent plusieurs analyses de certaines parties de l'image \footcite{redmon2016}.  

Si la détection d’objets prévue par \gls{yolo} peut être appliqué à des flux vidéo, nous souhaitons l'intégrer comme composant complémentaire au processus de segmentation précédemment décrit. De fait, il être pertinent de réutiliser l'image clé fixe, produite et sauvegardée par l'outil pyscenedetect, associée à chaque séquence, afin de la faire analyser par le modèle. Ceci permettrait de réduire le temps de calcul du processus et d'améliorer l'efficacité des prédictions, puisque le modèle reçoit un minimum d’informations.  

À intégrer ? “YOLO impose des contraintes spatiales strictes aux prédictions de boîtes englobantes, puisque chaque cellule de la grille ne prédit que deux boîtes et ne peut contenir qu’une seule classe. Cette contrainte spatiale limite le nombre d’objets proches que notre modèle peut prédire. Notre modèle éprouve des difficultés avec les petits objets qui apparaissent en groupes, comme les volées d’oiseaux. »\footcite{redmon2016} ("YOLO imposes strong spatial constraints on bounding box predictions since each grid cell only predicts two boxes and can only have one class. This spatial constraint limits the number of nearby objects that our model can predict. Our model struggles with small objects that appear in groups, such as flocks of birds." \footcite{redmon2016} 

Ainsi, l’application du modèle \gls{yolo} s’impose comme la suite logique du pipeline de segmentation détaillé précédemment, afin d’indexer les segments produits, via l’utilisation des \gls{imgcle}. Si le premier processus de segmentation produit des données temporelles associées à chaque séquence sous la forme de time codes, \gls{yolo} produit des données sémantiques sur ces mêmes séquences. À l'issue de l'analyse d'une \gls{imgcle}, le modèle \gls{yolo} est capable de produire trois types de données : des étiquettes sur les objets identifiés, un score de confiance qui indique l'indice de certitude concernant la prédiction, des coordonnées spatiales de localisation de l'objet détecté dans l'image \footcite{redmon2016}. Les étiquettes produites sur les objets identifiés, permettant de décrire les caractéristiques de l’image, pourront être indexées comme des mots-clés au sein du logiciel de gestion des collections. Ces mots-clés seront associés aux time-codes de chaque séquence détectée. Ainsi, chaque séquence contiendrait les données suivantes : des time codes qui l'identifient comme extrait, des étiquettes/mots-clés sur les objets qu'elle contient, et des données qui permettraient aux utilisateurs de vérifier la prédiction et de la valider. Le score de confiance produit par le modèle, permet en effet d'indiquer à l'utilisateur la certitude de l'étiquette associée à tel objet détecté. De même, les coordonnées spatiales de l’étiquette produite permettent à l’utilisateur de vérifier dans l'image si l'étiquette produite est associée au bon objet.  



\subsubsection{Entraînement}

à intégrer ? : %alternatives non supervisées : la supervision demande des données annotées (caron et al), l'apprentissage non supervisé (Unsupervised Learning) comme DeepCluster (Caron et al., 2018). Ce type de modèle regroupe les images ou frames similaires en "clusters" algorithmiques (via K-means) pour s'auto-entraîner sans intervention humaine (Caron et al., 2018). Si DeepCluster s'avère robuste face à des changements de distribution d'images (ce qui est prometteur pour des archives spécifiques), il présente deux limites majeures : il introduit des biais imprévisibles en remplaçant les labels par des métadonnées brutes (inacceptable pour la rigueur historienne) et a été conçu sur des bases gigantesques (ImageNet, 1.28 million d'images), le rendant difficilement transposable à des collections d'archives réduites (Caron et al., 2018 ; Norindr, 2023).

Pour les deux sous sections suivantes, nous prendront comme exemple d’application les archives britanniques en ligne, qui proposent une chaine de traitement pour la détection d’objet basée sur l'\gls{Apprentissage profond} et sur comment intégrer la supervision humaine dans la chaine de traitement et de vérification. 
%lien exact : voir doc word   

Toutefois, afin d’analyser le contenu des images, les \gls{cnn} subissent un apprentissage supervisé (Supervised Learning), c’est-à-dire un entraînement au cours duquel les algorithmes apprennent la détection d’objets à partir de données labellisées par l'humain \footcite{delua2024}. L'efficacité des modèles basés sur l'\gls{Apprentissage profond} dépend fondamentalement de nature et de la qualité de leur données d'entraînement. Yolov8 a fait l'objet d'un entraînement préalable en amont de sa mise à disposition, sur le jeu de données COCO, prévu notamment pour la détection d'objets, et composé de multiples objets du quotidien en tous genres \footcite{2023b}. Il est composé de 200 000 images annotées où les objets sont entourés et décrits. Il s'agit d'un jeu de données généraliste qui permet au modèle de détecter des objets communs tels qu'une voiture, un chat, un parapluie.  

Or, l’ancienneté et la qualité d’image du corpus d'archives audiovisuelles étudié, en fait un corpus d’images spécifique sur lequel l’algorithme doit être entraîné. Afin d'adapter le modèle à notre corpus et ainsi d'en améliorer l'efficacité, une phase de pré traitement est nécessaire sur les données du corpus, au cours de laquelle il est des données sont sélectionnée afin d’être annotées et testées \footcite{kasturi2024}. La supervision humaine est déterminante dans ce processus, particulièrement dans le cas d’analyse de corpus inédits sur lesquels l’algorithme reste à entraîner \footcite{kasturi2024}. Les british online archives implémentent par exemple l'adaptation du modèle en deux étapes (dans quel projet?? : %https://shura.shu.ac.uk/32356/3/2023_ukci_object_detection_in_archival_documents_with_hil.pdf )
, dans le cadre de la mise en place d’un pipeline de détection d’objets sur un corpus d’images d’archives. Certaines archives sont sélectionnées et annotées manuellement à l'aide de la plateforme \gls{labels} \footcite{kasturi2024}. 
%présernter un peu plus la plateforme label studio 
La méthode du transfert d’apprentissage (transfert learning) est ensuite appliquée, au cours de laquelle les données annotées servent d'entraînement de détection d'objets au modèle \footcite{kasturi2024}. Le transfert d'apprentissage permet notamment de \textit{fine tuner} un modèle existant, c’est-à-dire de l’affiner, afin de l'adapter au corpus analysé \footcite{bhargav2019}. Ainsi, la même méthode pourra être appliquée au corpus étudié. Le modèle existant, entraîné en amont sur des données génériques, pourra être affiné sur un jeu de données restreint annoté, plus spécifique \footcite{bhargav2019}. Ces deux étapes, l’annotation des données et l’affinage du modèle, constituent des étapes de prétraitement, en amont de l’analyse des données, afin d'améliorer l'efficacité de la détection d’objets \footcite{kasturi2024}. Concrètement, certaines des \gls{imgcle} produites par le processus de segmentation pourraient être sélectionnées afin d’être annotées. Ces annotations, dédiées à la détection d’objets, prennent la forme de détourages sur les objets ou caractéristiques que le modèle doit détecter, associées à des étiquettes \footcite{du2024}. L'outil \gls{labels} permet notamment la réalisation de ces annotations sur les images, à partir desquelles le modèle est entraîné à une détection spécifique. 

Cette étape de pré-traitement, avant le processus d'analyse, implique la définition des types de mots clés souhaités, pour l'indexation des séquences. En effet, les étiquettes associées aux objets au sein des images correspondent aux mots-clés qui seront prédits par le modèle lors de l’analyse. Dans le cadre du corpus étudié, il est pertinent d’extraire les mots-clés déjà indexés sur les \gls{sequence} au sein du logiciel. Ces mots-clés varient selon le contenu de l’image et peuvent, par exemple, être les suivants : affiche, foule, policier, gendarme, drapeau, etc. Ces mots-clés, tout comme le jeu de données COCO, sont assez génériques. De fait, le pré traitement est avant tout utile pour l’entraînement du modèle à pouvoir détecter des objets au sein d’images dont la qualité est compromise.   


\subsubsection{Interface de validation technique}

Afin de gérer les essais-erreurs de l'affinage du modèle sur les données du corpus, il est pertinent de mettre en place une interface de validation technique et d'amélioration du modèle. Cette interface permettrait d'une part de valider les résultats du modèle, et d’autre part, d'assurer la concordance entre les données produites par le modèle avec les données déjà existantes au sein du logiciel.  

En effet, suite à l'affinage du modèle, les archives britanniques en ligne préconisent des étapes supplémentaires de validation des données produites \footcite{kasturi2024}. Cette validation se structure en différentes étapes qui consistent notamment en l'analyse des données produites par le modèle, la modification et la correction de ces données, le potentiel réentraînement du modèle avec de nouveaux paramètres et/ou de nouvelles annotations, et enfin, l'évaluation des résultats sur les données de test \footcite{kasturi2024}. Cette validation peut s'effectuer sur la même interface que celle proposée pour l'annotation des données, label studio, qui permet la modification et la correction des annotations. Les modifications des annotations sont ainsi réinjectées dans le modèle afin d'améliorer ses performances \footcite{kasturi2024}. Est ce que c’est aussi simple ? 

Décrire plus précisément comment label studio interagit avec yolo : format de fichier, seuil de confiance 

Utiliser le score de confiance pour la validation : les étiquettes qui ot un score de confiance elevée sont valides, les autre sont vérifiées et re annotées.  

De plus, afin de pouvoir valider les données produites par le modèle, il faut pouvoir valider leur indexation au sein du logiciel, en adéquation avec les règles métier en place. Au sein de l’espace de travail S-Movies dédié au corpus étudié, l'indexation des mots-clés est contrôlée par le thésaurus "keyword". L'affinage et la correction des données produites par le modèle devra ainsi prendre en compte la concordance des mots-clés générés avec le thésaurus de mots clés contrôlés dans l'application. Or, le modèle Yolo est nativement pré entraîné sur un jeu de données généraliste duquel la majorité des prédictions seront puisées. Ainsi, afin d'éviter la génération de mots clés trop éloignés des termes du thésaurus dans l'application, il pourrait être pertinent d'établir un mapping, c'est-à-dire une table de correspondance entre les mots clés que le modèle peut potentiellement générer, et leur forme contrôlée dans le thésaurus, pour aider la validation technique des données produites par le modèle. Nous avons effectivement déjà mis ce procédé en place au cours de notre stage dans le cadre de la mise en place de l'intelligence artificielle pour d'autres tâches, notamment la mise en place d'une recherche en langage naturel. (Mettre une note de bas de page pour expliquer). Pour cette tâche tout à fait différente que celle que nous décrivons dans ce mémoire, nous avons effectué une table de correspondance entre les champs contrôlés pouvant être requêtés dans le logiciel et leur définition en langage naturel, préparant ainsi la gestion de synonymes et de termes employés par l'utilisateur potentiellement pour requêter ces champs. Sur le même principe, nous pourrions mettre en place une table de correspondance afin d'assurer la validation des données produites et d'assurer \textit{in fine} leur indexation au sein de l'application.  

La finalité est de pouvoir permettre à l'utilisateur de pouvoir recherche ces mots clés et d'accéder aux segments concernés. Transition vers la prochaine sous partie avec ça  

Nécessité de l’active learning : pour que le modèle soit développé et cohérent, il faudrait qu’il soit entrainé et réentrainé en permanence, à mesure qu’il reçoit des données client et que les données sont validées ou non  

citation supplémentaire sur les détails des corrections d'entraînement :  

"The evaluation employs two distinct test datasets, namely test-1 and test2. One of these datasets will be utilised for the purpose of rectifying the predictions, while the other dataset will be exclusively utilised for evaluating the model’s performance across varying levels of HIL. Upon the user’s modification of the bounding boxes on one of the test dataset, the corresponding corrected images will replace some random images in the initial training dataset which then becomes a new training dataset. A small learning rate of 0.00005 (in our case) is employed to retrain the model using the new training dataset. To retain previously learned information while incorporating the user-provided data, we utilize transfer-learning techniques that load the pre-trained weights of the model from HIL-0%. This approach allows us to make gradual adjustments to the model’s weights, ensuring the assimilation of the new data without compromising the existing knowledge" \footcite{kasturi2024} 


Transition avec la partie suivante : la validation finale est laissée à l'utilisateur métier : interface de validation livrée au client, une fois que les entraînements et affinages sont concluants. 




généralement, un algorithme de machine learning comprend trois étapes : data pre processing, data modelling, process optimisation. Dans le pre processing : humain requis pour transformer des données non structurées en données structurées et annotées \footcite{kasturi2024}

 










